\chapter{Introducción}
\label{cap:capitulo1}
\setcounter{page}{1}




\section{Contextualización: la contaminación digital}


La digitalización ha transformado la sociedad y la economía, pero su avance ha generado un impacto ambiental significativo y a menudo subestimado, conocido como \textbf{Contaminación Digital}. El sector de las Tecnologías de la Información y Comunicación (TIC) es un consumidor energético masivo y un emisor relevante de Gases de Efecto Invernadero (GEI).


Este impacto se complejiza por la distinción entre el \textbf{Carbono Operacional} (consumo de energía durante el uso) y el creciente \textbf{Carbono Incorporado} (emisiones generadas en la fabricación y disposición del \textit{hardware}). La tendencia hacia el uso intensivo de tecnologías como la Inteligencia Artificial y el \textit{Machine Learning} amplifica esta presión energética, posicionando al \textit{software} y los algoritmos como nuevos focos de mitigación.


\section{El vacío curricular en sostenibilidad digital}


El problema central radica en que, a pesar de esta huella ecológica documentada, la sostenibilidad digital no se ha integrado de manera sistemática en la formación práctica de los futuros profesionales de la Informática. Los currículos actuales suelen centrarse en la funcionalidad y la eficiencia económica, relegando la eficiencia energética y la conciencia ambiental a un plano secundario.

Este trabajo parte de la necesidad de dotar a los estudiantes de Informática de herramientas teóricas y prácticas para mitigar este impacto. Para ello, se propone trascender la mera concienciación e incorporar los principios de la \textbf{Educación para el Desarrollo Sostenible (EDS)}, cuyo objetivo es desarrollar un cambio conductual que integre conocimiento, actitudes y comportamiento sostenible.
